\chapter[Open Problems]{Open Problems}\label{ch:sec:open}
The deterministic realization problem is completely characterized in this book.  Several natural extensions remain open.
\paragraph{Part I--Weighted Incidence Structures.}
\begin{itemize}
\item \textbf{Probabilistic encoders.}  Allowing randomized encoding strategies may relax the integrality condition on the multiplicity ratios.  A characterization of probabilistically realizable WIS would be a natural next step.
\item \textbf{Infinite message spaces.}  Extending the WIS framework to countable or continuous message spaces requires measure-theoretic formulations of the quotient geometry and the universal optimality defect.
\item \textbf{Continuous weighted incidence structures.}  When the incidence graph is replaced by a continuous bipartite structure (e.g.\ a kernel), the balance equation \(B\beta = Kw\) becomes an integral equation.  The existence of solutions and their interpretation as encoding strategies is unexplored.
\item \textbf{Optimization of the universal optimality defect.}  The defect \(\delta(Y)\) defines a cost function on the space of WIS weights.  Minimizing \(\delta\) over admissible weights, subject to integrality constraints, may yield approximate encoders for systems where exact universal optimality is unattainable.
\item \textbf{Approximation theory for irrational weights.}  When the induced multiplicity ratios are irrational, no exact finite encoder exists.  Rational approximation of the weights and the resulting trade-off between approximation quality and key length merit investigation.
\item \textbf{Categorical formulations.}  The WIS construction, the gauge quotient, and the Hilbert-space structure may admit a unified description in terms of category-theoretic fibrations, providing a bridge to the theory of information algebras.
\end{itemize}
\paragraph{Part II--Hilbert Geometry.}
\begin{itemize}
\item \textbf{Spectral geometry of the posterior Laplacian.}  The reduced Hessian derived in Chapter~8 is a graph Laplacian on the incidence graph.  The relationship between its spectrum and the combinatorial properties of the original encoding system (key length, message entropy, equivocation) is largely unexplored.
\item \textbf{Non-quadratic variational principles.}  The least-squares formulation of Chapter~8 is forced by the logarithmic linearization.  Higher-order Taylor expansions of the WIS factorization may give rise to non-quadratic variational problems whose analysis requires methods from non-linear functional analysis.
\end{itemize}
\paragraph{Part III--Posterior Manifold Geometry.}
\begin{itemize}
\item \textbf{Global curvature.}  The pullback metric defined in Part~III is flat because it is pulled back from Euclidean space.  Replacing the Euclidean target with a curved space (e.g.\ the Fisher--Rao metric on the codomain) would yield a non-flat geometry whose curvature encodes additional information about the encoding system.
\item \textbf{Boundary behaviour of entropy bounds.}  The Lipschitz and quadratic bounds of Part~IV break down near the boundary of the probability simplex.  A refined analysis for distributions with small but non-zero probabilities---the regime relevant to practical side-channel analysis---remains open.
\end{itemize}
\paragraph{Part IV--Information Functionals and Security.}
\begin{itemize}
\item \textbf{Tightness of the Lipschitz constants.}  The constants \(L_S\) and \(C_S\) in the universal bounds of Part~IV are obtained from supremum bounds on the WIS gradient.  Determining the optimal (smallest) constants for specific functionals---Shannon entropy, R\'enyi entropy, min-entropy---would sharpen the estimates and clarify the limits of the geometric approach.
\item \textbf{Beyond smooth functionals.}  The piecewise-smooth theory of Chapter~17 covers rank-based security measures.  Extending the analysis to functionals with more complex singularities (e.g.\ the guessing entropy for non-uniform distributions) would broaden the applicability of the framework.
\end{itemize}
\paragraph{Part V--Variational Structure.}
\begin{itemize}
\item \textbf{Multiplicity of critical points.}  Theorem~49 (Three-Level Structure) shows that interior critical points of the uniform observer divergence collapse to at most three distinct log-likelihood values.  Whether analogous collapse phenomena occur for non-uniform reference measures is an open question.
\item \textbf{Convergence of gradient flows.}  The variational formulation of the UOD suggests a natural gradient flow on the posterior manifold.  Establishing convergence guarantees for this flow would connect the geometric theory to practical optimization algorithms for approximate universal optimality.
\end{itemize}
\paragraph{Cross-cutting directions.}
\begin{itemize}
\item \textbf{Quantum WIS.}  Generalizing weighted incidence structures to non-commutative probability spaces would yield a quantum analogue of the universal optimality defect, with potential applications to quantum key distribution and quantum side-channel analysis.
\item \textbf{Computational complexity.}  The cycle-factorization criterion (Theorem~4.1), the balance-equation test for realizability (Theorem~5.2), and the computation of the universal optimality defect all have polynomial-time algorithms.  A systematic complexity-theoretic study of the associated decision and optimization problems is lacking.
\end{itemize}
\endinput
